\documentclass[a4paper,11pt]{ltjsarticle}
\usepackage{amsmath,amssymb,amsthm,mathtools,hyperref}
\usepackage{luatexja-fontspec}
\setmainjfont{HaranoAjiMincho-Regular}
\newtheorem{theorem}{定理}[section]
\newtheorem{lemma}[theorem]{補題}
\newtheorem{proposition}[theorem]{命題}
\newtheorem{corollary}[theorem]{系}
\newtheorem{remark}[theorem]{注意}
\renewcommand{\proofname}{証明}
\newcommand{\Perf}{\operatorname{Perf}}
\newcommand{\RHom}{\operatorname{RHom}}
\newcommand{\thick}{\operatorname{thick}}
\newcommand{\otL}{\mathbin{\otimes^{\mathbf L}}}
\newcommand{\Spec}{\operatorname{Spec}}
\title{非連結dg代数上のperfectnessの反例と完備基底での限界}
\author{非可換cubic threefold研究・ソース復元用ノート}
\date{2026年11月10日・v1}
\begin{document}
\maketitle
\begin{abstract}
\texttt{notes/pdf-only/finite\_smoothness\_notes.pdf}（2026年11月1日版）
第23節（pp.~26--31）の数学的内容をTeXで独立に再構成する。
特殊ファイバーとすべてのArtin切断がperfectであっても、
完備基底上の相対dg加群がperfectとは限らないことを証明する。
この反例の特殊ファイバーの代数自体はsmoothではなく、相対対角の
smoothness持ち上げに対する反例ではない。
本稿は元の34ページのPDF全体を再生成するソースではない。
\end{abstract}

\section{規約と問題}
すべてcohomological gradingとし、微分の次数を$+1$とする。
$\Perf(\Lambda)$はunbounded derived category $D(\Lambda)$のcompact objects、
同値に有限coneとretractで$\Lambda$から生成される対象を表す。
基底環上のperfectnessとdg加群としてのperfectnessを混同しない。
次の含意は一般に偽である：
\[
\Lambda,M\in\Perf(R),\qquad
M\otL_Rk\in\Perf(\Lambda\otL_Rk)
\quad\Longrightarrow\quad M\in\Perf(\Lambda).
\tag{DG4.1}
\]

\section{特殊ファイバーだけperfectな二生成子の例}
\begin{theorem}[特殊ファイバーとすべてのArtin切断でperfect]\label{thm:two}
$R=k[[h]]$, $K=k((h))$とし
\[
\Lambda=R\oplus R\varepsilon\oplus R\eta,\qquad
|\varepsilon|=|\eta|=1,\qquad
(\varepsilon,\eta)^2=0,\quad d_\Lambda=0.
\]
右dg加群$M=e\Lambda$は$|e|=0$, $d_M(e)=he\varepsilon$とする。
このとき
\begin{enumerate}
\item $\Lambda,M$は$R$上の有界有限自由複体である。
\item $M_k=M\otL_Rk\simeq\Lambda_k$は$\Lambda_k$-perfectである。
\item $M_K=M\otL_RK$は$\Lambda_K$-perfectでない。
\item すべての$n\ge1$で$M_n=M\otL_RR/(h^n)$は$\Lambda_n$-perfectである。
\item $M$は$\Lambda$-perfectでない。
\end{enumerate}
\end{theorem}
\begin{proof}
$(\varepsilon,\eta)^2=0$からLeibniz則で$d_M^2=0$となる。
基礎の$R$複体は
\[
0\longrightarrow R\xrightarrow{(h,0)}
R\varepsilon\oplus R\eta\longrightarrow0
\]
であるから1が従い、$h=0$で微分が零となるので2が従う。
$h$を可逆にすると$H^0(M_K)=0$, $H^1(M_K)=Ke\eta$である。
augmentation $\Lambda_K\to K$を通して$K$を加群とみれば
\[
K[-1]\xrightarrow{\simeq}M_K,\qquad1\longmapsto e\eta
\]
は$\Lambda_K$-linear quasi-isomorphismである。

$K$が$\Lambda_K$-perfectでないことを示す。
$\varepsilon,\eta$のすべての有限語$w$（空語を含む）に
次数0の自由生成子$e_w$を対応させる左semifree加群$Q$をとり、
\[
d(e_\varnothing)=0,\qquad d(e_{aw})=ae_w
\quad (a\in\{\varepsilon,\eta\})
\]
とする。radical square zeroゆえ$d^2=0$。
語長filtrationはsemifree filtrationであり、
非空語とradicalの元が対応して消えるので$Q\to K$は
quasi-isomorphismである。augmentationの右$K$とtensorすると
全微分が零になるから
\[
K\otL_{\Lambda_K}K\simeq
\bigoplus_{w\in\{\varepsilon,\eta\}^{*}}Ke_w
\]
は無限次元のcohomologyを持つ。
もし$K\in\Perf(\Lambda_K)$なら
$-\otL_{\Lambda_K}K$はこれを$K$-perfectに送るため矛盾する。
従って3。5はperfectnessが導来基底変更で保存されることから従う。
4は系\ref{cor:artin}から従う。
\end{proof}

\begin{corollary}[fiberwise perfect locusの非開性]
$R=k[h]$上で同じ式を定義すると
\[
\{p\in\Spec R\mid M\otL_R\kappa(p)\in
\Perf(\Lambda\otL_R\kappa(p))\}=V(h).
\]
\end{corollary}
\begin{proof}
$h=0$では自由加群であり、$h\ne0$の剰余体では
定理\ref{thm:two}のaugmentationを用いた議論が成立する。
\end{proof}
\begin{remark}
本例は$H^{>0}(\Lambda_k)\ne0$であり、
$H^{>0}=0$というconnectivity仮定を伴う定理への反例ではない。
さらに$\Lambda_k$自体はsmoothではない。
従って$E_k$ smoothを仮定する相対対角$E$の
perfectnessの反例を構成したものではない。
\end{remark}

\section{Artin切断とthick detection}
\begin{lemma}[thick detection]\label{lem:thick}
$R$を可換環、$S$を可換$R$代数とし
$R\in\thick_{D(R)}(S)$と仮定する。
$R$-dg代数$\Gamma$と$P\in D(\Gamma)$について
\[
P\otL_RS\in\Perf(\Gamma\otL_RS)
\quad\Longrightarrow\quad P\in\Perf(\Gamma)
\]
が成り立つ。connectivityは不要である。
\end{lemma}
\begin{proof}
$X\in D(R)$に対し函手
$N\mapsto\RHom_\Gamma(P,N\otL_RX)$が
任意の直和を保つという条件を考える。
この条件を満たす$X$はthick部分圏をなす。
$X=S$ではextension--restriction adjunctionによって
\[
\RHom_\Gamma(P,N\otL_RS)
\simeq \RHom_{\Gamma\otL_RS}
(P\otL_RS,N\otL_RS)
\]
となり、仮定されたcompactnessと基底変更の
直和保存性から右辺は直和を保つ。
よって$S$はこのthick部分圏に入り、仮定から$R$も入る。
従って$P$はcompactである。
\end{proof}
\begin{corollary}\label{cor:artin}
$R_n=k[h]/(h^n)$上のdg moduleのperfectnessは
$k$への導来基底変更で検出される。特に定理\ref{thm:two}の
すべての$M_n$は$\Lambda_n$-perfectである。
\end{corollary}
\begin{proof}
$R_n$は$k$を組成因子に持つ有限filtrationを持ち、
$R_n\in\thick_{D(R_n)}(k)$である。
補題\ref{lem:thick}を適用する。
\end{proof}
したがって$M\simeq\mathbf R\varprojlim M_n$かつ各$M_n$が
perfectでも$M$自身の$\Lambda$-perfectnessは従わない。

\section{全ファイバーがperfectでも相対perfectでない一生成子の例}
\begin{proposition}
$R=k[[h]]$または$k[h]_{(h)}$とし
\[
\Gamma=R\oplus R\varepsilon,\quad
|\varepsilon|=1,\quad\varepsilon^2=0,\quad d_\Gamma=0,
\qquad P=e\Gamma,\quad d_P(e)=he\varepsilon.
\]
$P$は$R$-perfectで、$P_k\simeq\Gamma_k$および
$P_K\simeq0$はいずれもperfectだが
$P\notin\Perf(\Gamma)$である。
\end{proposition}
\begin{proof}
基礎の$R$複体は$R\xrightarrow{h}R$。
仮に$P$が$\Gamma$-perfectならaugmentation
$\Gamma\to R$により$N=P\otL_\Gamma R\in\Perf(R)$。
導来基底変更では$N\otL_Rk\simeq k$だが
$N\otL_RK\simeq0$である。
局所環上のperfect complexは有界有限自由複体で表されるため
Euler characteristicは各項のrankの交代和であり、
special/genericで一致すべきところ$1\ne0$となる。矛盾。
\end{proof}

\section{完備基底でなお成立するHom有限性}
\begin{proposition}\label{prop:hom}
$R=k[[h]]$, $\Gamma$を$R$-dg代数、
$P,N\in D(\Gamma)$とする。
\[
P\otL_Rk\in\Perf(\Gamma\otL_Rk),\qquad
N\in\Perf(R)
\]
ならば$\RHom_\Gamma(P,N)\in\Perf(R)$である。
\end{proposition}
\begin{proof}
$C=\RHom_\Gamma(P,N)$とする。
$k$は$R$-perfectで$-\otL_Rk$は$h$倍写像のconeだから、
$P$のcompactnessを仮定せず
\[
C\otL_Rk\simeq
\RHom_{\Gamma\otL_Rk}(P\otL_Rk,N\otL_Rk)
\]
であり、右辺は$k$-perfect。
$N$は$R$-perfectゆえderived $h$-completeで、
$\RHom_R(R[1/h],N)=0$。
tensor--Hom adjunctionから
\[
\RHom_R(R[1/h],C)
\simeq\RHom_\Gamma(P,\RHom_R(R[1/h],N))=0.
\]
従って$C$はderived $h$-completeであり
$C\simeq\mathbf R\varprojlim_n(C\otL_RR/(h^n))$。
可換完備環上のperfect complexのeffectivity
（Stacks Project, \href{https://stacks.math.columbia.edu/tag/09AW}{Tag 09AW}）
から$C\in\Perf(R)$を得る。
\end{proof}
\begin{remark}
上の結論は$P$の$\Gamma$-compactnessを意味しない。
$R[1/h]$への一般の$\RHom$基底変更は、
$k$への有限coneでの基底変更とは異なる。
\end{remark}
\begin{corollary}
$E$が$R$-perfectな相対dg代数として構成済みで
$E_k$がsmoothなら、
\[
\RHom_{E^{\mathrm{op}}\otL_RE}(E,E)\in\Perf(R).
\]
これはHochschild cochainの有限性であり
$E$の相対smoothnessを意味しない。
\end{corollary}

\section{非可換cubicの正則性研究への適用限界}
Artin切断で対角がperfectであることから
完備基底上の対角の一様な有限cell/retract模型を
無条件に得ることはできない。
上の反例は一般の加群補題の使用を否定するが、
特定の非可換cubicのtailsがsmoothでないという
結論は与えない。中心局所化と有限行列因子化など、
対角の追加構造を利用した証明と区別する。

\paragraph{対応原本と未照合の範囲}
原本は\texttt{notes/pdf-only/finite\_smoothness\_notes.pdf}
の第23節（pp.~26--31）であり、
上の命題・証明について今回TeX化した。
原本34ページのうち他の部分はまだ対応確認できていないため、
このファイルだけを根拠に原本PDFをGit追跡解除してはならない。
\end{document}